\documentclass{situationist-preprint}

\title{Procreation and Consent}
\msid{ETH-20-b}
\msfamily{ETH-20}
\msdiscipline{Normative and Applied Ethics}
\mstype{Companion manuscript}
\mskeywords{procreation; consent; Shiffrin; pure benefits; antinatalism; parental responsibility}

\begin{document}
\maketitle

\begin{abstract}
ETH-20-a shows that Benatar's asymmetry fails, so that existence is not a harm to a person whose life contains more good than bad. A different objection to procreation does not depend on the asymmetry: creating a child imposes serious risks of harm on someone who has not consented and cannot consent, for the sake of benefits that prevent no greater harm. We argue that this objection does not show procreation to be wrong. The requirement of consent protects the authority of an existing will over its own life, and there is no such will to bypass; the standard that governs decisions for those who cannot consent is their interests, which procreation can serve. What the objection does establish is a responsibility of parents for the harms their children suffer.
\end{abstract}

\section{The objection}

\citet{Shiffrin1999} argues that there is an asymmetry between harms and benefits in the morality of imposition. We may impose a harm on a person without her consent to prevent a greater harm to her, but not to confer a \emph{pure benefit}, one that does not avert harm. Her illustration is a benefactor who, wishing to enrich the inhabitants of an island, drops valuable cubes of gold on it, knowing that some will be injured by them. Coming into existence, on her view, confers pure benefits and imposes serious harms, and the child's consent cannot be obtained. Shiffrin concludes that procreation is morally hazardous and that procreators bear responsibility for the harms their children suffer. \citet{Benatar2006} enlists considerations of this kind in support of the stronger conclusion that procreation is wrong. That stronger conclusion is our target.

\section{The argument}

\begin{proposition}\label{p:noauthority}
The wrong of imposing a harm without consent, where consent is unobtainable, is the wrong of overriding the authority of an existing person over her own life; no such authority is overridden by creating a person.
\end{proposition}

\begin{proof}[Argument]
In the gold-cube case the islanders exist, have their own plans and are entitled to decide which risks to run for which goods. The benefactor wrongs them because he decides for them a question that is theirs to decide, and the wrong does not depend on the balance of benefit and harm, which may be favorable. What makes consent required is the existence of a will entitled to decide. Before conception there is no will whose decision the procreator preempts, and no plan of life into which the risks are thrust; the person created acquires authority over her life only by being created, and from that point her parents may not override it.
\end{proof}

Where consent cannot be had from someone who will have interests, the governing standard is substituted decision: one acts in accordance with the interests of the person affected, as we do for infants and for those who cannot consent \citep[cf.][]{Harman2004}. ETH-20-a shows that, on the impersonal standard that parity requires, existence is better than never existing for a life containing more good than bad (its Proposition~3.5 and Theorem~3.6). A procreator who has good reason to expect such a life, and who will act to secure it, decides in accordance with the interests of the person created.

\begin{proposition}\label{p:residue}
Procreation that meets the substituted-decision standard is not wrong on the ground of lack of consent; it leaves the procreator responsible for the harms the child suffers, in the sense of owing her care, mitigation and, where possible, remedy.
\end{proposition}

\begin{proof}[Argument]
By Proposition~\ref{p:noauthority}, the ground of the consent requirement is absent; by the substituted-decision standard and ETH-20-a, the decision can be made in the child's interests. What remains of Shiffrin's asymmetry is that the procreator has, by her own choice, exposed someone to harms she did not choose. That is a ground of responsibility for those harms, and the duties of parents, which outlast childhood, discharge it.
\end{proof}

\begin{objection}
The substituted-decision standard is meant for persons who already exist and cannot consent. Applying it to the not-yet-existent is as arbitrary as the asymmetry ETH-20-a rejects.
\end{objection}

\begin{reply}
The standard applies to anyone whose interests a decision will affect and whose consent cannot be had. A future person has interests in the relevant sense: there will be someone whose life goes better or worse because of the decision, and that is all the standard requires. What the not-yet-existent lacks is not interests but a present will, and it is the absence of the will, not of interests, that removes the consent requirement.
\end{reply}

\begin{objection}
The risks are not trivial: some lives contain more bad than good, and no procreator can exclude this.
\end{objection}

\begin{reply}
Substituted decisions are judged by the reasonableness of the risk given what is at stake, as in the medical treatment of children. A procreator with good reason to expect a life worth living takes a reasonable risk, and Proposition~\ref{p:residue} assigns her the duty to respond if it materializes. A procreator with good reason to expect a life containing more bad than good does not, and the objection then succeeds; but it succeeds on the ground of the child's interests, not of consent.
\end{reply}

\section{Conclusion}

The consent objection does not show that having children is wrong. It shows that procreators answer for what they impose. Together with ETH-20-a, this leaves no general ground for antinatalism in the child's own case.

\begin{thebibliography}{9}\small
\bibitem[Benatar(2006)]{Benatar2006} Benatar, D. (2006). \emph{Better Never to Have Been: The Harm of Coming into Existence}. Oxford University Press.
\bibitem[Harman(2004)]{Harman2004} Harman, E. (2004). Can we harm and benefit in creating? \emph{Philosophical Perspectives}, 18.
\bibitem[Shiffrin(1999)]{Shiffrin1999} Shiffrin, S.~V. (1999). Wrongful life, procreative responsibility, and the significance of harm. \emph{Legal Theory}, 5(2), 117--148.
\end{thebibliography}

\end{document}
